\newpage
\Article{73}
What we have set out in Arts.~69---72 is, strictly speaking, independent of the assumption that $m$ is a primary number. From now on, however, we shall at least suppose that $a$ is odd and $b$ even, and further that $a,b$ and $a-b$ are positive numbers. First of all, we must establish the bounds on the values of $x$ in the collection $G$.

\SourceLeaf{154}%
Putting $ay-bx=\eta,(a+b)x-(a-b)y=\zeta,p-2ax-2by=\theta$, the criterion for numbers $x+yi$ to belong to the collection $G$ consists in the three conditions that $\eta,\zeta,\theta$ be positive numbers. Since $px=(a-b)\eta+a\zeta,p(a-2x)=a\theta+2b\eta$, it is evident that $x$ and $2a-x$ must be positive numbers, or that $x$ must equal one of the numbers $1,2,3\ldots\frac12(a-1)$. Further, since $(a-b)\theta=2b\zeta+p(a-b-2x)$, it is clear that, as long as $x$ is less than $\frac12(a-b)$, the second condition (requiring $\zeta$ to be positive) already implies the third (that $\theta$ must be positive). Conversely, whenever $x$ is greater than $\frac12(a-b)$, the second condition is already contained in the third. Thus, for the values $x$ takes in $1,2,3\ldots\frac12(a-b-1)$, we need only ensure that $\eta$ and $\zeta$ are positive, or that $y$ is greater than $\frac{bx}{a}$ and less than $\frac{(a+b)x}{a-b}$. For such a given value of $x$, the total count of numbers $x+yi$ is therefore
\[
\left[\frac{(a+b)x}{a-b}\right]-\left[\frac{bx}{a}\right]
\]
if we use the brackets in the same sense in which we have already used them repeatedly elsewhere (compare \emph{Theorematis arithm. dem. nova}, Art.~4, and \emph{Theorematis fund. in doctr. de residuis quadr.} etc., \emph{Algorithm. nov.}, Art.~3). For the values $x$ takes in $\frac12(a-b+1),\frac12(a-b+3)\dotsfour\frac12(a-1)$, on the other hand, it suffices to give $\eta$ and $\theta$ positive values, or to make $y$ greater than $\frac{bx}{a}$ and less than $\frac{p-2ax}{2b}$, or $\frac12b+\frac{aa-2ax}{2b}$. Thus, for such a given value of $x$, the total count of numbers $x+yi$ is
\[
\left[\frac12b+\frac{aa-2ax}{2b}\right]-\left[\frac{bx}{a}\right]
\]
Hence we conclude that the count of numbers in the collection $G$ is
\[
g=\Sigma\left[\frac{(a+b)x}{a-b}\right]+\Sigma\left[\frac12b+\frac{aa-2ax}{2b}\right]-\Sigma\left[\frac{bx}{a}\right]
\]
where the first sum extends over all integer values of $x$ from $1$ to $\frac12(a-b-1)$, the second from $\frac12(a-b+1)$ to $\frac12(a-1)$, and the third from $1$ to $\frac12(a-1)$.

\Needspace{14\baselineskip}
If we use the symbol $\varphi$ in the same sense as in the place cited (\emph{Theorematis fund.} etc., \emph{Algor. nov.}, Art.~3), namely so that
\[
\varphi(t,u)=\left[\frac{u}{t}\right]+\left[\frac{2u}{t}\right]+\left[\frac{3u}{t}\right]\dotsfour+\left[\frac{t'u}{t}\right]
\]
where $t,u$ denote any positive numbers and $t'$ the number $[\frac12t]$, then the first term becomes $=\varphi(a-b,a+b)$, and the third $=-\varphi(a,b)$; the second, however, becomes
\SourceLeaf{155}%
\[
=\frac14bb+\Sigma\left[\frac{aa-2ax}{2b}\right]
\]
But, writing the terms in reverse order, we have
\[
\begin{aligned}
\Sigma\left[\frac{aa-2ax}{2b}\right]&=\left[\frac{a}{2b}\right]+\left[\frac{3a}{2b}\right]+\left[\frac{5a}{5b}\right]+\dotsfour+\left[\frac{(b-1)a}{2b}\right]\\
&=\varphi(2b,a)-\varphi(b,a)
\end{aligned}
\]
Our formula therefore takes the following form:
\[
g=\varphi(a-b,a+b)+\varphi(2b,a)-\varphi(a,b)-\varphi(b,a)+\frac14bb
\]

Let us first consider the term $\varphi(a-b,a+b)$, which immediately changes into $\varphi(a-b,2b)+1+2+3+\text{ etc. }+\frac12(a-b-1)$, or into
\[
\varphi(a-b,2b)+\frac18((a-b)^2-1)
\]

Next, since the general theorem gives $\varphi(t,u)+\varphi(u,t)=[\frac12t]\pdot[\frac12u]$ when $t,u$ are relatively prime positive integers, we have
\[
\varphi(a-b,2b)=\frac12b(a-b-1)-\varphi(2b,a-b)
\]
and consequently
\[
\varphi(a-b,a+b)=\frac18(aa+2ab-3bb-4b-1)-\varphi(2b,a-b)
\]

Let us arrange the parts of $\varphi(2b,a-b)$ as follows:
\[
\begin{aligned}
&\left[\frac{a-b}{2b}\right]+\left[\frac{3(a-b)}{2b}\right]+\left[\frac{5(a-b)}{2b}\right]+\text{ etc. }+\left[\frac{(b-1)(a-b)}{2b}\right]\\
&+\left[\frac{a-b}{b}\right]+\left[\frac{2(a-b)}{b}\right]+\left[\frac{3(a-b)}{b}\right]+\text{ etc. }+\left[\frac{\frac12b(a-b)}{b}\right]
\end{aligned}
\]
The second series evidently becomes
\[
\begin{aligned}
&=\varphi(b,a-b)=\varphi(b,a)-1-2-3-\text{ etc. }-\frac12b\\
&=\varphi(b,a)-\frac18(bb+2b)
\end{aligned}
\]
We write the first series, with the order of its terms reversed, as follows:
\[
\begin{aligned}
&\left[\frac12(a+1-b)-\frac{a}{2b}\right]+\left[\frac12(a+3-b)-\frac{3a}{2b}\right]\\
&+\left[\frac12(a+5-b)-\frac{5a}{2b}\right]+\text{etc.}+\left[\frac12(a-1)-\frac{(b-1)a}{2b}\right]
\end{aligned}
\]
This expression. Since, with $t$ denoting an integer and $u$ a fraction, we generally have $[t-u]=t-1-[u]$, it changes into the following:
\[
\begin{aligned}
&\frac18b(2a-4-b)-\left[\frac{a}{2b}\right]-\left[\frac{3a}{2b}\right]-\left[\frac{5a}{2b}\right]-\text{ etc. }-\left[\frac{(b-1)a}{2b}\right]\\
&=\frac18b(2a-4-b)-\varphi(2b,a)+\varphi(b,a)
\end{aligned}
\]

\SourceLeaf{156}%
Hence
\[
\varphi(2b,a-b)=2\varphi(b,a)-\varphi(2b,a)+\frac14b(a-3-b)
\]
and therefore
\[
\varphi(a-b,a+b)=\varphi(2b,a)-2\varphi(b,a)+\frac18(aa-bb+2b-1)
\]
Substituting this value in the formula for $g$ given above, and also substituting $\varphi(a,b)+\varphi(b,a)=\frac14b(a-1)$, we obtain
\[
g=2\varphi(2b,a)-2\varphi(b,a)+\frac18(aa-2ab+bb+4b-1)
\]

\newpage
\Article{74}
Precisely similar reasoning deals with the case where $a$, $b$ remain positive but $a-b$ is negative, or $b-a$ positive. The equations $p(a-2x)=2b\eta+a\theta$, $p(b-a+2x)=2b\zeta+(b-a)\theta$ show that $\frac12a-x$ and $x+\frac12(b-a)$ must be positive, and hence that $x$ must equal one of the numbers $-\frac12(b-a-1),-\frac12(b-a-3),-\frac12(b-a-5)\dotsfour+\frac12(a-1)$. Further, the equation $px+(b-a)\eta=a\zeta$ shows that, for negative values of $x$, the condition requiring $\eta$ to be positive is already contained in the condition requiring $\zeta$ to be positive, whereas the reverse holds whenever $x$ is given a positive value. Thus the values of $y$, for a fixed negative value of $x$, must lie between $\frac{(a+b)x}{a-b}$ and $\frac{p-2ax}{2b}$, whereas for a positive value of $x$ they must lie between $\frac{bx}{a}$ and $\frac{p-2ax}{2b}$. For $x=0$, these bounds are evidently $0$ and $\frac{p-2ax}{2b}$, with the value $y=0$ itself excluded. Hence
\[
g=-\Sigma\left[\frac{(a+b)x}{a-b}\right]+\Sigma\left[\frac12b+\frac{aa-2ax}{2b}\right]-\Sigma\left[\frac{bx}{a}\right]
\]
where the first sum extends over all negative values of. $x$ from $-1$ to $-\frac12(b-a-1)$; the second over all values of $x$ from $-\frac12(b-a-1)$ to $\frac12(a-1)$; and the third over all positive values of $x$ from $+1$ to $\frac12(a-1)$. In this way the first sum gives $-\varphi(b-a,b+a)$, the second, as in the preceding article, $\frac14bb+\varphi(2b,a)-\varphi(b,a)$, and finally the third $-\varphi(a,b)$; thus we have
\[
g=-\varphi(b-a,b+a)+\varphi(2b,a)-\varphi(b,a)-\varphi(a,b)+\frac14bb
\]
Now, in the same way as in the preceding article, we obtain
\SourceLeaf{157}%
\[
\begin{aligned}
\varphi(b-a,b+a)&=\varphi(b-a,2b)-\frac18((b-a)^2-1)\\
&=\frac18(3bb-2ab-aa-4b+1)-\varphi(2b,b-a)
\end{aligned}
\]
as well as
\[
\varphi(2b,b-a)=\varphi(2b,a)-2\varphi(b,a)+\frac14b(b-1-a)
\]
and consequently
\[
\varphi(b-a,b+a)=2\varphi(b,a)-\varphi(2b,a)+\frac18(bb-aa-2b+1)
\]
and finally
\[
g=2\varphi(2b,a)-2\varphi(b,a)+\frac18(aa-2ab+bb+4b-1)
\]
It is therefore proved that the same formula for $g$ holds whether $a-b$ is positive or negative, provided that $a$, $b$ are positive.

\Article{75}
To obtain a further reduction, we put
\[
\begin{aligned}
L&=\left[\frac{a}{2b}\right]+\left[\frac{2a}{2b}\right]+\left[\frac{3a}{2b}\right]+\text{etc.}+\left[\frac{\frac12ba}{2b}\right]\\
M&=\left[\frac{(\frac12b+1)a}{2b}\right]+\left[\frac{(\frac12b+2)a}{2b}\right]+\left[\frac{(\frac12b+3)a}{2b}\right]+\text{etc.}+\left[\frac{ba}{2b}\right]\\
N&=\left[\frac{a+b}{2b}\right]+\left[\frac{2a+b}{2b}\right]+\left[\frac{3a+b}{2b}\right]+\text{etc.}+\left[\frac{\frac12ba+b}{2b}\right]
\end{aligned}
\]
It is easy to see that $[u]+[u+\frac12]=[2u]$ holds in general, whatever real quantity $u$ denotes. Thus $L+N=\varphi(b,a)$, and since evidently $L+M=\varphi(2b,a)$, we have
\[
\varphi(2b,a)-\varphi(b,a)=M-N
\]

Further, it is evident that the sum of the first term of the series $N$ and the next-to-last term of the series $M$, namely $\left[\frac{a+b}{2b}\right]+\left[\frac{(b-1)a}{2b}\right]$, is $=\frac12(a-1)$, and that the same sum is formed by the second term of the series $N$ and the third-to-last term of the series $M$, and so on. Since the last term of the series $M$ is also $=\frac12(a-1)$, while the last term of the series $N$ is $=\left[\frac{a+2}{4}\right]=\frac14(a\mp1)$, with the upper or lower sign applying according as $a$ is of the form $4n+1$ or $4n-1$, we have
\[
M+N=\frac14(a-1)b+\frac14(a\mp1)
\]
and therefore
\[
\varphi(2b,a)-\varphi(b,a)=\frac14(a-1)b+\frac14(a\mp1)-2N
\]

\SourceLeaf{158}%
Thus the formula for $g$ found in Arts.~73 and 74 becomes
\[
g=\frac18((a+b)^2-1)+2n-4N
\]
putting $a\mp1=4n$, where $n$ will be an integer. Since this gives $1=16nn-8an+aa$, this formula can also be written as follows:
\[
g=\frac18(-aa+2ab+bb+1)+4(\frac12(a+1)n-nn-N)
\]
Since $g$ is the character of the number $-1-i$ for the modulus $a+bi$, this character is therefore $\equiv\frac18(-aa+2ab+bb+1)\pmod4$. This is precisely the theorem found above (Art.~64) by examining examples, and the theorems concerning the characters of the numbers $1+i$, $1-i$, $-1+i$ follow from it at once. Thus these four theorems have now been rigorously proved for the case where $a$ and $b$ are positive.

\Needspace{16\baselineskip}
\Article{76}
If $a$ remains positive and $b$ is negative, put $b=-b'$, so that $b'$ is positive. It is now proved that, for the modulus $a+b'i$, the character of the number $-1-i$ is $\equiv\frac18(-aa+2ab'+b'b'+1)\pmod4$. Hence the character of the number $-1+i$ for the modulus $a-b'i$, by the theorem stated in Art.~62, is $\equiv\frac18(aa-2ab'-b'b'-1)$, that is, the character of the number $-1+i$ for the modulus $a+bi$ becomes $\equiv\frac18(aa+2ab-bb-1)$. This is precisely the theorem stated in Art.~64, from which the other three concerning the characters of the numbers $1+i$, $1-i$, $-1-i$ follow at once. Thus these theorems are also proved for the case where $b$ is negative, that is, for all cases where $a$ is positive.

Finally, if $a$ is negative, put $a=-a'$, $b=-b'$. By what has already been proved, the character of the number $1+i$ with respect to the modulus $a'+b'i$ is $\equiv\frac18(-a'a'+2a'b'-3b'b'+1)\pmod4$, and it makes no difference whether we take the number $a'+b'i$ or its opposite $-a'-b'i$ as the modulus. Thus the character of the number $1+i$ with respect to the modulus $a+bi$ is evidently $\equiv\frac18(-aa+2ab-3bb+1)$, and the same applies to the characters of the numbers $1-i$, $-1+i$, $-1-i$.

We therefore conclude that the proof of the theorems concerning the characters of the numbers $1+i$, $1-i$, $-1+i$, $-1-i$ (Arts.~63. 64) is no longer subject to any restriction.
\BookRule{.28\linewidth}
